\section{基本制度}
\begin{enumerate}
  \item \zd{合议制}是重点。实践中程序空转，为了提高司法效率，扩大了独任制的适用范围。
  \item 司法民主。
        \begin{enumerate}
          \item 参审。让陪审员和法官一起组成合议庭，一起对事实的认定、法律的适用问题发表意见，进行评议。
          \item \zd{陪审}。事实认定和法律认定分开。我国。
        \end{enumerate}
  \item 审判组织。\zd{独任制}、\zd{合议制}（\bdy 普通程序）、审委会。少数服从多数原则。\zd{考各种审判组织的适用}。
        \begin{enumerate}
          \item 一审。事实认定。
                \begin{enumerate}
                  \item 简易程序 \yjt 独任制。
                  \item 普通程序 \yi 合议制。\er \zd{独任制}。基层法院审理 + 基本事实清楚、权利关系明确。
                  \item \zd{一审可以陪}。人民陪审员只参加适用\zd{一审程序审理的诉讼程序}。
                  \item \zd{一审只要是中院，就不能是独任制}。
                \end{enumerate}
          \item 二审。吸收当事人对裁判结果不满的功能，定纷止争。事实和法律适用。
                \begin{enumerate}
                  \item 原则合议制。
                  \item \zd{独任制}。新增。
                        \begin{enumerate}
                          \item 级别。中院。
                          \item 一审是简易程序 \hz 不服一审裁定。
                          \item 事实清楚、权利义务关系明确。
                          \item \zd{经双方当事人同意}。
                          \item 发现独任不适合的，应裁定转为合议庭审理。
                          \item 提异议的。\yi 成立 \yjt 转为合议庭。\er 不成立 \yjt 驳回。
                        \end{enumerate}
                \end{enumerate}
          \item 再审合议庭组成具体看是几审。\yi 一审 \yjt 一审。\er 二审 \hz 提审 \yjt 二审。
          \item 特别程序独任制；选民资格和重大疑难案件 \yjt 合议制。
          \item 督促程序 = 支付令 \yjt 独任制。
          \item 公示催告程序。公示催告阶段 \yjt 独任制；除权判决阶段 \yjt 合议制。
        \end{enumerate}
  \item 回避。\zd{用决定，不是裁定}。
        \begin{enumerate}
          \item 适用对象。\zd{期待保持中立}。
                \begin{enumerate}
                  \item 与案件有\zd{利害关系}。
                  \item \zd{证人，具有不可替代性，虽然需要中立，但不适用回避制度}。
                \end{enumerate}
          \item 回避的法定原因。很多都是常识问题。
                \begin{enumerate}
                  \item 其他利害关系=能说得出关系，看是否影响公正审理，做题看\zd{有没有修饰性的词语}。
                  \item 有不正当的行为。\al{请客送礼、给当事人推荐律师、违规会见当事人}（\zd{非工作时间、工作地点}）。
                  \item 参与过前续审判。例外。发回重审 \yjt 一审 \yjt 二审，\zd{原二审合议庭成员}除外。
                \end{enumerate}
          \item 考试只考申请回避。看题目出现了几次回避，是不是都是指同一个回避，不是的话每个回避单独处理。
                \begin{enumerate}
                  \item 口头、书面都可以。
                  \item \zd{开始审理时}提出，最迟\zd{法庭辩论终结前}。
                  \item \zd{审判人员}的回避由院长决定。\zd{判决书有名字的} \yjt 院长，没名的 \yjt 审判长。
                  \item 换人\zd{程序不需要重来}；被申请人\zd{暂停工作到决定作出前}，采取紧急措施的除外。
                  \item 提出 \yjt 驳回 \yjt \zd{申请人可以复议一次}，期间\zd{不停止}原决定的执行。
                \end{enumerate}
        \end{enumerate}
  \item 公开审判制度。
        \begin{enumerate}
          \item 法定不公开。国家秘密、个人隐私、其他案件。
          \item 申请不公开。离婚诉讼、商业秘密。\zd{一方}申请。
          \item \zd{评议一定不公开}，但宣判一律要\zd{公开宣判}。
        \end{enumerate}
  \item 两审终审制度。例外 \yjt 一审终审（含判决、裁定）。
        \begin{enumerate}
          \item 最高法作出的一审判决。
          \item 特别程序、督促程序、公示催告程序。
          \item 小额诉讼程序的实体判决、驳回起诉、管辖权异议裁定。
          \item 一般的裁定书。例外：不予受理、驳回起诉、管辖权异议。
          \item 调解书。
        \end{enumerate}
\end{enumerate}
